% Generated by scripts/2026-10-21/extract_floats.py from paper/main_en.tex; edit the full manuscript, not this file.
\begin{table}[htbp]
\centering\small
\caption{Confirmatory hold-out (episodes 250--499; 3B reader, 75\% depth, linear).}\label{tab:confirm}
\begin{tabularx}{\linewidth}{@{}Lccc@{}}
\toprule
 & Step 0 & Step 1 & Step 2 \\
\midrule
Probe AUROC (95\% CI) & 0.645 (0.547, 0.736) & 0.716 (0.619, 0.810) & 0.729 (0.633, 0.815) \\
Rung 5 AUROC & 0.607 & 0.681 & 0.758 \\
Probe $-$ rung 5 (95\% CI) & $+0.039$ ($-0.030$, $+0.100$) & $+0.035$ ($-0.030$, $+0.096$) & $-0.029$ ($-0.087$, $+0.029$) \\
Decision (margin 0.05, set post hoc) & non-inferior & non-inferior & inconclusive \\
\bottomrule
\end{tabularx}
\par\smallskip\raggedright\footnotesize
$n = 250$ episodes (205 failures). Training used all 500 episodes; evaluation used only the new 250. Margin sensitivity, read from the same intervals (lower bounds $-0.030$, $-0.030$, $-0.087$; upper bounds all above zero): with a margin of 0.025 no step would be non-inferior, and with 0.10 all three would be. Because the margin was set post hoc, the differences and intervals, not the labels, carry the evidence.
\end{table}
